% !TEX root = ../discretemath.tex

\section{Существование равновесия Нэша через теорему Какутани}


\subsection{Теорема существования}

\begin{Theorem}{О существовании равновесия Нэша}{}
	В любой биматричной игре существует хотя бы одно равновесие Нэша (в смешанных стратегиях).
\end{Theorem}

\subsection{Теорема Какутани}

\begin{Theorem}{Какутани о неподвижной точке}{}
	Пусть $S\subset\mathbb{R}^n$ — непустое компактное выпуклое множество, $F\colon S\to 2^S$ — многозначное отображение такое, что:
	\begin{enumerate}
		\item $F(x)\ne\varnothing$ для всех $x\in S$;
		\item $F(x)$ выпукло для всех $x\in S$;
		\item $F$ \emph{замкнуто}: если $x_k\to x$, $y_k\to y$ и $y_k\in F(x_k)$, то $y\in F(x)$.
	\end{enumerate}
	Тогда существует неподвижная точка: $\exists\,x^*\in S$ с $x^*\in F(x^*)$.
\end{Theorem}

\begin{Remark}{}{}
	Теорема Какутани обобщает теорему Брауэра (непрерывное отображение замкнутого шара в себя имеет неподвижную точку) на случай многозначных замкнутых отображений с выпуклыми значениями.
\end{Remark}

\begin{figure}[H]
	\centering
	\begin{tikzpicture}[font=\small, >=Stealth, x=1.2cm, y=1.2cm]
		\draw[->] (-0.2,0) -- (3.4,0) node[right]{$x$};
		\draw[->] (0,-0.2) -- (0,3.4) node[above]{$F(x)$};
		\draw (3,0)--(3,3)--(0,3);
		% диагональ y=x
		\draw[dashed] (0,0) -- (3,3) node[above left]{$y=x$};
		% многозначное отображение со скачком (вертикальный отрезок = выпуклое значение)
		\draw[very thick, blue] (0,0.6) -- (1.5,1.2);
		\draw[very thick, blue] (1.5,1.2) -- (1.5,2.4);
		\draw[very thick, blue] (1.5,2.4) -- (3,2.7);
		% неподвижная точка на пересечении вертикального куска с диагональю
		\filldraw[red] (1.5,1.5) circle (2.4pt);
		\draw[dashed] (1.5,0)--(1.5,1.5); \node[below] at (1.5,0) {$x^*$};
		\node[red!70!black, align=center] at (2.5,1.0) {\scriptsize $x^*\in F(x^*)$};
	\end{tikzpicture}
	\caption{Идея теоремы Какутани (одномерная схема): «толстый» график замкнутого многозначного отображения с выпуклыми значениями (вертикальный отрезок — значение $F$ в точке скачка) обязан пересечь диагональ $y=x$. Точка пересечения — неподвижная: $x^*\in F(x^*)$.}
\end{figure}

\subsection{Доказательство существования равновесия Нэша}

\begin{proof}
	Рассмотрим компактное выпуклое множество всех пар смешанных стратегий
	\[
	S=\Bigl\{(\vec p,\vec q):\ p_i\ge 0,\ \sum_i p_i=1,\ q_j\ge 0,\ \sum_j q_j=1\Bigr\}\subset\mathbb{R}^{m+n}.
	\]
	Определим многозначное отображение «наилучших ответов» $F\colon S\to 2^S$:
	\[
	F(\vec p,\vec q)=\Bigl\{(\vec p^{\,\prime},\vec q^{\,\prime})\in S:\
	M_1(\vec p^{\,\prime},\vec q)=\max_{\vec x}M_1(\vec x,\vec q),\
	M_2(\vec p,\vec q^{\,\prime})=\max_{\vec y}M_2(\vec p,\vec y)\Bigr\}.
	\]
	Проверим условия теоремы Какутани.
	\begin{enumerate}
		\item \textbf{Непустота.} Максимум непрерывной функции на компакте $S$ достигается, поэтому $F(\vec p,\vec q)\ne\varnothing$.
		\item \textbf{Выпуклость.} Функции $M_1(\,\cdot\,,\vec q)$ и $M_2(\vec p,\,\cdot\,)$ линейны по своему аргументу, поэтому множество максимизаторов выпукло; значит, и $F(\vec p,\vec q)$ выпукло.
		\item \textbf{Замкнутость.} Если $(\vec p^{(k)},\vec q^{(k)})\to(\vec p,\vec q)$ и $(\vec p^{\prime(k)},\vec q^{\prime(k)})\to(\vec p^{\,\prime},\vec q^{\,\prime})$ с $(\vec p^{\prime(k)},\vec q^{\prime(k)})\in F(\vec p^{(k)},\vec q^{(k)})$, то по непрерывности $M_1,M_2$ предельные значения тоже максимальны, то есть $(\vec p^{\,\prime},\vec q^{\,\prime})\in F(\vec p,\vec q)$.
	\end{enumerate}
	По теореме Какутани существует $(\vec p^{\,*},\vec q^{\,*})\in F(\vec p^{\,*},\vec q^{\,*})$, что в точности означает
	\[
	M_1(\vec p^{\,*},\vec q^{\,*})=\max_{\vec x}M_1(\vec x,\vec q^{\,*}),\qquad
	M_2(\vec p^{\,*},\vec q^{\,*})=\max_{\vec y}M_2(\vec p^{\,*},\vec y),
	\]
	то есть условие равновесия Нэша: ни одному игроку не выгодно отклоняться в одностороннем порядке.
\end{proof}

